\section{Discussion and limitations}
\label{sec:discussion}

In homogeneous Case~H, the rotated characteristic-length tensor changes the gradient-stiffness contribution according to the wave-vector direction relative to its principal axes. The resulting branch shifts are unequal: rotating the tensor can separate neighboring bands along $\Gamma$--$X$ while modes at other wave vectors remain in the same frequency interval. Under the volume-equivalent scaling, changing the aspect ratio lengthens one principal characteristic length while shortening the other, so the two bounding band edges need not move together. Their frequency difference can therefore vary non-monotonically. A positive $\Delta_{GX}$ is therefore a reduced-zone separation along $\Gamma$--$X$ only, with modes elsewhere in the Brillouin zone consistent with the negative complete-gap measure $\Delta_{\mathrm{complete}}\le -0.3235$; without material contrast it is direction-selective spectral separation, not full-zone attenuation of an incident wave. The 42-point map compares evaluated configurations rather than identifying a continuous optimum.

The group-velocity response follows from the same direction-dependent dispersion. Because $\bmv_g=\nabla_{\bmk}\omega$ is normal to a local constant-frequency contour, changes in contour slope alter the angle between the energy-flow direction and the wave vector. Rotating the length tensor reorients these local contour normals and shifts the steering lobes. On the sampled acoustic branch at $\bar{k}=0.5$, the peak deviation is nearly unchanged while its principal direction moves. Orientation therefore redirects the energy flow without materially increasing its peak angular deviation, within the branch, ring, and orientations sampled here.

Case~C adds periodic contrast in both shear modulus and density. The resulting Bloch bands separate Bands~3 and~4 over the sampled Brillouin zone, unlike the path-specific branch separation in homogeneous Case~H. This difference is consistent with the positive complete-gap measure in Case~C and the absence of a complete gap in Case~H. The Case-C gap decreases from $2.5732$ on the $4\times4$ mesh to $1.9179$ on the $64\times64$ mesh and remains unconverged. Because the cases use different material and discretization settings, the comparison is interpretive rather than a controlled parameter sweep; the coarse-mesh Case-C width is not a continuum estimate.

The homogeneous mesh study gives an empirical rate $p=4.17$ over the $4^2$, $8^2$, and $16^2$ levels. The $32^2$ result is resolution-limited and excluded from the fit, so this rate is not a theoretical-order claim. All eight tabulated numerical consistency checks meet their stated criteria; the energy-flux/group-velocity identity, evaluated on closed-form Case-H harmonic fields, holds to $6.91\times10^{-10}$.

The classical AlN/BaTiO$_3$ calculation reproduces the Rytov transfer-matrix relation to $2.22\times10^{-16}$, although the comparison with the published B1 curves stays visual because tabulated source values are unavailable. The two gradient-elastic bilayer cases are reproductions rather than quantitative external validations, since their source parameter conventions are unresolved; no external-validation claim is made for those benchmarks.

For the transverse branch in ~\ref{app:asymptotics}, gradient stiffness raises the phase speed at high wavenumber when micro-inertia is absent. With micro-inertia, the $k^2$ contribution in the numerator of $v_p^2$ is balanced asymptotically by the $k^2$ inertial term in its denominator, leaving a finite limit set by the ratio of the gradient-stiffness and micro-inertia lengths. This is a mechanism-specific result for the branch and constitutive model studied, not a universal bound. More broadly, the model is linear, local, isothermal, and two-dimensional under plane-strain assumptions; it excludes nonlinear gradient effects, dissipation, piezoelectric coupling, and three-dimensional polarization modes.
